\begin{frame}{프로젝트 채점 기준 — Block B 3주가 수렴하는 곳}
  \small
  {\footnotesize 8.1 rubric이 그대로 이어져 오면서 ``Ch4$\sim$7 개념''이
  ``Ch9$\sim$15 개념''으로 교체된 것이 전부:}
  \vskip0.4em
  \begin{tabular}{@{}llp{0.42\textwidth}@{}}
    \toprule
    항목 & 비중 & 좋은 답이란 \\
    \midrule
    문제 정의와 환경 & 15\% & 상태·행동·보상이 명확(물리적 단위까지,
    Ch13), 환경 선택이 Ch9$\sim$15의 개념 검증 목적에 근거 \\
    알고리즘 적절성 & 20\% & 환경 특성(이산/연속, 희소성, 모델 유무)에
    알고리즘이 맞고, 근거가 Ch9$\sim$15 개념으로 설명 \\
    \textbf{실험 프로토콜의 정직성} & \textbf{25\%} & 시드 $\geq 3$,
    학습 곡선과 최종 평가(탐험 끔) 분리, 베이스라인(무작위 또는 사람이
    설계한 PD 제어 등), $\varepsilon\cdot\gamma\cdot\lambda\cdot$클립
    $\epsilon$ 감수성 실험이 보고서에서 추적 가능 \\
    수식적 정당화 & 20\% & 배운 개념 $\geq 1$개로 ``왜''를 설명(11.3
    클리핑, 12.1 복합 오차, 13.1 도메인 무작위화 $\ldots$)했고,
    관측 숫자와 일치 \\
    \textbf{결과의 정직성과 반성} & \textbf{20\%} & 불안정했던 시드·구간
    을 숨기지 않고 원인을 Ch9$\sim$15 개념으로 진단(16.1 좋은/나쁜 예
    기준) \\
    \bottomrule
  \end{tabular}
  \vskip0.5em
  \textbf{정직성 항목 = 합계 45\%} — 목적은 ``좋은 정책''이 아니라
  ``그 \textbf{성능 숫자를 정직하게 얻는 절차}를 실제로 수행했음을
  증명''. greedy 평가 $-1100$이 나와도 프로토콜이 정직하면 높은 점수,
  ``수렴했다''고 주장해도 시드 1개였다면 크게 감점 — 8.1의 ``\textbf{한
  번의 실행은 아무것도 말해주지 않는다}''가 연속 제어·신경망 무대에서
  다시 쓰인 것. 16.2의 5항목이 이 표의 열 \emph{그대로}인 이유도
  여기: \textbf{리뷰어와 채점자는 같은 눈으로 본 것}.
\end{frame}
